\subsection{** APS‑TSLCA-SECTION-003 **}\label{apstslca-section-003}

\subsection{** Three-Squared-Lattice Cognitive Architecture **}\label{three-squared-lattice-cognitive-architecture}

\subsection{** Symbiotic Universal Xessability Standards **}\label{symbiotic-universal-xessability-standards}

\subsection{** Aurphyx Primordial Standard **}\label{aurphyx-primordial-standard}

\subsection{** Aurphyx LLC **}\label{aurphyx-llc}

\subsection{** SAGES \textbar{} Proprietary \textbar{} Pro-Existence **}\label{sages-proprietary-pro-existence}

\subsection{** Accessibility = Xessability **}\label{accessibility-xessability}

\subsection{** Version 3.69 **}\label{version-3.69}

\subsection{3. The Three aXes}\label{the-three-axes}

The Three-Squared-Lattice Cognitive Architecture is built on three orthonormal cognitive aXes. Each aXis governs a distinct dimension of intelligence, and together they form the basis vectors of the cognitive lattice. This section defines each aXis formally, establishes its operational domain, and clarifies its role within the broader cognitive field.

\begin{center}\rule{0.5\linewidth}{0.5pt}\end{center}

\subsubsection{3.1 Somatic Intelligence aXis (SIX)}\label{sensorimotor-integration-axis-six}

SIX governs the system's coupling with the world. It is responsible for perception, accessibility, embodiment, and environmental attunement. In the lattice, SIX defines the \textbf{perceptual axis}, ensuring that cognition is grounded in real sensory data and contextual signals.

SIX includes: - multimodal perception (visual, auditory, tactile, proprioceptive) - accessibility transformations (SAIL, VAP, environmental overlays) - sensorimotor integration and feedback loops - embodied context (location, motion, orientation) - environmental salience and affordance detection

SIX ensures that the cognitive field remains anchored to the external world and that all semantic and identity operations remain grounded in lived context.

\begin{center}\rule{0.5\linewidth}{0.5pt}\end{center}

\subsubsection{3.2 Systemic Coherence aXis (SCX)}\label{systemic-coherence-axis-scx}

SCX governs semantics, reasoning, invariants, and system‑level coherence. It defines the \textbf{semantic axis} of the lattice, providing the structures that allow information to be interpreted, transformed, and maintained consistently across time.

SCX includes: - semantic parsing and transmutation (Fuxyez, FUTE) - invariant enforcement and coherence maintenance - systemic reasoning and global context propagation - rule‑based and field‑based interpretation - integration with AuraFS as the topological semantic substrate

SCX ensures that cognition remains logically consistent, semantically stable, and capable of supporting complex reasoning across distributed or embodied systems.

\begin{center}\rule{0.5\linewidth}{0.5pt}\end{center}

\subsubsection{3.3 Identity Coherence aXis (ICX)}\label{soul-identity-axis-icx}

ICX governs identity continuity, provenance, self‑consistency, and ethical grounding. It defines the \textbf{identity axis} of the lattice, ensuring that the system maintains a coherent sense of self across time, transformations, and contexts.

ICX includes: - identity anchoring (SoulShot, BlissID) - provenance tracking and GuardHash lattice binding - ethical invariants and SAGES enforcement - memory continuity and self‑referential stability - value alignment and contextual self‑modulation

ICX ensures that cognition remains stable across time, that actions remain consistent with identity and values, and that all transformations preserve continuity of self.

\begin{center}\rule{0.5\linewidth}{0.5pt}\end{center}

\subsubsection{3.4 Orthogonality of the aXes}\label{orthogonality-of-the-axes}

The three aXes are orthogonal in the cognitive vector space:

\[
\mathbf{s}_i \cdot \mathbf{s}_j = 0 \quad (i \neq j).
\]

This orthogonality ensures: - perception does not collapse into semantics\\
- semantics does not overwrite identity\\
- identity does not distort perception

Each aXis is independent yet combinable, forming the basis for the nine‑cell lattice.

\begin{center}\rule{0.5\linewidth}{0.5pt}\end{center}

\subsubsection{3.5 Interdependence and Mutual Reinforcement}\label{interdependence-and-mutual-reinforcement}

Although orthogonal, the aXiss are interdependent. Cognition emerges from their interactions: - SIX grounds SCX in reality.\\
- SCX interprets SIX into meaning.\\
- ICX stabilizes both through continuity and ethics.

This triadic structure is the minimal configuration capable of supporting stable, reversible, identity‑preserving cognition.

\begin{center}\rule{0.5\linewidth}{0.5pt}\end{center}

\subsubsection{3.6 aXes as Field Generators}\label{axes-as-field-generators}

Each aXis generates a cognitive sub‑field:

\[
\Phi_{\text{SIX}}, \quad \Phi_{\text{SCX}}, \quad \Phi_{\text{ICX}}.
\]

The full cognitive field is the superposition:

\[
\Phi = \Phi_{\text{SIX}} + \Phi_{\text{SCX}} + \Phi_{\text{ICX}}.
\]

This field evolves according to invariants enforced by SAGES, ensuring coherence, stability, and ethical alignment.

\begin{center}\rule{0.5\linewidth}{0.5pt}\end{center}

\subsubsection{3.7 aXes as Basis Vectors of the Lattice}\label{axes-as-basis-vectors-of-the-lattice}

The aXiss form the basis vectors:

\[
\mathbf{s}_1 = \text{SIX}, \quad \mathbf{s}_2 = \text{SCX}, \quad \mathbf{s}_3 = \text{ICX}.
\]

Their tensor products generate the nine cognitive cells:

\[
\mathbf{s}_i \otimes \mathbf{s}_j.
\]

These cells define the operational modes of cognition, from raw perception to identity‑modulated semantics.

\begin{center}\rule{0.5\linewidth}{0.5pt}\end{center}

\subsubsection{3.8 aXes as Interfaces}\label{axes-as-interfaces}

Each aXis also functions as an interface layer: - SIX interfaces with sensors, actuators, and accessibility systems.\\
- SCX interfaces with semantic engines, reasoning modules, and topological computation.\\
- ICX interfaces with identity systems, provenance chains, and governance layers.

This interface structure ensures that the cognitive lattice remains substrate‑agnostic and universally implementable.
